
\documentclass{article}
\usepackage{colortbl}
\usepackage{makecell}
\usepackage{multirow}
\usepackage{supertabular}

\begin{document}

\newcounter{utterance}

\centering \large Interaction Transcript for game `dond', experiment `coop\_en', episode 1 with claude{-}sonnet{-}4{-}5{-}20250929{-}t1.0.
\vspace{24pt}

{ \footnotesize  \setcounter{utterance}{1}
\setlength{\tabcolsep}{0pt}
\begin{supertabular}{c@{$\;$}|p{.15\linewidth}@{}p{.15\linewidth}p{.15\linewidth}p{.15\linewidth}p{.15\linewidth}p{.15\linewidth}}
    \# & $\;$A & \multicolumn{4}{c}{Game Master} & $\;\:$B\\
    \hline

    \theutterance \stepcounter{utterance}  
    & & \multicolumn{4}{p{0.6\linewidth}}{
        \cellcolor[rgb]{0.9,0.9,0.9}{
            \makecell[{{p{\linewidth}}}]{
                \texttt{\tiny{[P1$\langle$GM]}}
                \texttt{You are playing a cooperative negotiation game in which you have to agree on how to divide a set of items among you and another player.} \\
\\ 
\texttt{Rules:} \\
\texttt{(a) You and the other player are given a set of items. Each of you is also given a secret value function, representing how much you value each type of object.} \\
\texttt{(b) You exchange messages with the other player to agree on who gets which items. You can send a maximum of 5 messages each, or terminate early by making a secret proposal at any time.} \\
\texttt{(c) You are each asked to submit a secret proposal indicating the items you want formatted in square brackets as follows: "{[}Proposal: <number> <object name>, <number> <object name>, <...>{]}"} \\
\texttt{(d) If your proposals are complementary, i.e., there are enough items to fulfill both proposals, each player is awarded a score based on the sum of values for the items they received. Otherwise, both of you get zero points.} \\
\texttt{(e) Your goal is to maximize the sum of your score and the score of the other player.} \\
\\ 
\texttt{Let us start.} \\
\\ 
\texttt{The set of available items is:} \\
\\ 
\texttt{2 bells, 1 doll, 2 threads} \\
\\ 
\texttt{Your secret value function is:} \\
\\ 
\texttt{bell: 1, doll: 2, thread: 3} \\
\\ 
\texttt{IMPORTANT: Your messages, unless it is the secret proposal, are directly transmitted to the other player, so do not include any response to the rules or text announcing your message. To make a secret proposal, use the indicated format. Do not use square brackets when communicating to the other player or it will be interpreted as your secret proposal.} \\
            }
        }
    }
    & \\ \\

    \theutterance \stepcounter{utterance}  
    & \multicolumn{4}{p{0.6\linewidth}}{
        \cellcolor[rgb]{0.9,0.9,0.9}{
            \makecell[{{p{\linewidth}}}]{
                \texttt{\tiny{[P1$\rangle$GM]}}
                \texttt{Hello! Let's work together to maximize our combined score. To do this efficiently, we should figure out which items each of us values more. } \\
\\ 
\texttt{I'd like to understand your preferences {-} could you share which items you value most? For my part, I have different values for each item type, so there's likely a way to divide things that benefits us both.} \\
            }
        }
    }
    & & \\ \\

    \theutterance \stepcounter{utterance}  
    & & & \multicolumn{2}{p{0.3\linewidth}}{
        \cellcolor[rgb]{0.9,0.9,0.9}{
            \makecell[{{p{\linewidth}}}]{
                \texttt{\tiny{[GM$|$GM]}}
                \texttt{continue} \\
            }
        }
    }
    & & \\ \\

    \theutterance \stepcounter{utterance}  
    & & & \multicolumn{2}{p{0.3\linewidth}}{
        \cellcolor[rgb]{0.9,0.9,0.9}{
            \makecell[{{p{\linewidth}}}]{
                \texttt{\tiny{[GM$|$GM]}}
                \texttt{Hello! Let's work together to maximize our combined score. To do this efficiently, we should figure out which items each of us values more. } \\
\\ 
\texttt{I'd like to understand your preferences {-} could you share which items you value most? For my part, I have different values for each item type, so there's likely a way to divide things that benefits us both.} \\
            }
        }
    }
    & & \\ \\

    \theutterance \stepcounter{utterance}  
    & & \multicolumn{4}{p{0.6\linewidth}}{
        \cellcolor[rgb]{0.9,0.9,0.9}{
            \makecell[{{p{\linewidth}}}]{
                \texttt{\tiny{[GM$\rangle$P2]}}
                \texttt{You are playing a cooperative negotiation game in which you have to agree on how to divide a set of items among you and another player.} \\
\\ 
\texttt{Rules:} \\
\texttt{(a) You and the other player are given a set of items. Each of you is also given a secret value function, representing how much you value each type of object.} \\
\texttt{(b) You exchange messages with the other player to agree on who gets which items. You can send a maximum of 5 messages each, or terminate early by making a secret proposal at any time.} \\
\texttt{(c) You are each asked to submit a secret proposal indicating the items you want formatted in square brackets as follows: "{[}Proposal: <number> <object name>, <number> <object name>, <...>{]}"} \\
\texttt{(d) If your proposals are complementary, i.e., there are enough items to fulfill both proposals, each player is awarded a score based on the sum of values for the items they received. Otherwise, both of you get zero points.} \\
\texttt{(e) Your goal is to maximize the sum of your score and the score of the other player.} \\
\\ 
\texttt{Let us start.} \\
\\ 
\texttt{The set of available items is:} \\
\\ 
\texttt{2 bells, 1 doll, 2 threads} \\
\\ 
\texttt{Your secret value function is:} \\
\\ 
\texttt{bell: 0, doll: 2, thread: 4} \\
\\ 
\texttt{IMPORTANT: Your messages, unless it is the secret proposal, are directly transmitted to the other player, so do not include any response to the rules or text announcing your message. To make a secret proposal, use the indicated format. Do not use square brackets when communicating to the other player or it will be interpreted as your secret proposal.} \\
\\ 
\\ 
\texttt{Hello! Let's work together to maximize our combined score. To do this efficiently, we should figure out which items each of us values more. } \\
\\ 
\texttt{I'd like to understand your preferences {-} could you share which items you value most? For my part, I have different values for each item type, so there's likely a way to divide things that benefits us both.} \\
            }
        }
    }
    & \\ \\

    \theutterance \stepcounter{utterance}  
    & & & \multicolumn{4}{p{0.6\linewidth}}{
        \cellcolor[rgb]{0.9,0.9,0.9}{
            \makecell[{{p{\linewidth}}}]{
                \texttt{\tiny{[GM$\langle$P2]}}
                \texttt{Hello! Great approach. I'm happy to share my preferences openly to find the best allocation.} \\
\\ 
\texttt{I value the items in this order: threads are most valuable to me, then the doll, and I don't value the bells at all.} \\
\\ 
\texttt{So ideally, I'd benefit most from getting the threads and the doll, while you could have both bells. But of course, this depends on your values {-} if you also value threads or the doll highly, we'll need to find a different split.} \\
\\ 
\texttt{What are your preferences? Which items matter most to you?} \\
            }
        }
    }
     \\ \\

    \theutterance \stepcounter{utterance}  
    & & & \multicolumn{2}{p{0.3\linewidth}}{
        \cellcolor[rgb]{0.9,0.9,0.9}{
            \makecell[{{p{\linewidth}}}]{
                \texttt{\tiny{[GM$|$GM]}}
                \texttt{continue} \\
            }
        }
    }
    & & \\ \\

    \theutterance \stepcounter{utterance}  
    & & & \multicolumn{2}{p{0.3\linewidth}}{
        \cellcolor[rgb]{0.9,0.9,0.9}{
            \makecell[{{p{\linewidth}}}]{
                \texttt{\tiny{[GM$|$GM]}}
                \texttt{Hello! Great approach. I'm happy to share my preferences openly to find the best allocation.} \\
\\ 
\texttt{I value the items in this order: threads are most valuable to me, then the doll, and I don't value the bells at all.} \\
\\ 
\texttt{So ideally, I'd benefit most from getting the threads and the doll, while you could have both bells. But of course, this depends on your values {-} if you also value threads or the doll highly, we'll need to find a different split.} \\
\\ 
\texttt{What are your preferences? Which items matter most to you?} \\
            }
        }
    }
    & & \\ \\

    \theutterance \stepcounter{utterance}  
    & & \multicolumn{4}{p{0.6\linewidth}}{
        \cellcolor[rgb]{0.9,0.9,0.9}{
            \makecell[{{p{\linewidth}}}]{
                \texttt{\tiny{[P1$\langle$GM]}}
                \texttt{Hello! Great approach. I'm happy to share my preferences openly to find the best allocation.} \\
\\ 
\texttt{I value the items in this order: threads are most valuable to me, then the doll, and I don't value the bells at all.} \\
\\ 
\texttt{So ideally, I'd benefit most from getting the threads and the doll, while you could have both bells. But of course, this depends on your values {-} if you also value threads or the doll highly, we'll need to find a different split.} \\
\\ 
\texttt{What are your preferences? Which items matter most to you?} \\
            }
        }
    }
    & \\ \\

    \theutterance \stepcounter{utterance}  
    & \multicolumn{4}{p{0.6\linewidth}}{
        \cellcolor[rgb]{0.9,0.9,0.9}{
            \makecell[{{p{\linewidth}}}]{
                \texttt{\tiny{[P1$\rangle$GM]}}
                \texttt{Thank you for being so open! That's very helpful.} \\
\\ 
\texttt{My preferences are actually similar to yours {-} I also value threads most, then the doll, then bells least. So we both want the threads most!} \\
\\ 
\texttt{Given this, we need to find the split that maximizes our combined total. Since you value bells at zero and I do get some value from them, the bells should definitely go to me.} \\
\\ 
\texttt{For the threads and doll {-} since we both value threads highly, the question is who values them MORE relative to the doll. If you value threads significantly more than the doll compared to my ratio, then you should get the threads. } \\
\\ 
\texttt{Based on your description that threads are "most valuable" and the doll is second, while for me the difference isn't as extreme, I think the optimal split is probably:} \\
\texttt{{-} You get: 2 threads and 1 doll} \\
\texttt{{-} I get: 2 bells} \\
\\ 
\texttt{This way you get your high{-}value items and I get something rather than nothing. Does this seem fair to you for maximizing our combined score?} \\
            }
        }
    }
    & & \\ \\

    \theutterance \stepcounter{utterance}  
    & & & \multicolumn{2}{p{0.3\linewidth}}{
        \cellcolor[rgb]{0.9,0.9,0.9}{
            \makecell[{{p{\linewidth}}}]{
                \texttt{\tiny{[GM$|$GM]}}
                \texttt{continue} \\
            }
        }
    }
    & & \\ \\

    \theutterance \stepcounter{utterance}  
    & & & \multicolumn{2}{p{0.3\linewidth}}{
        \cellcolor[rgb]{0.9,0.9,0.9}{
            \makecell[{{p{\linewidth}}}]{
                \texttt{\tiny{[GM$|$GM]}}
                \texttt{Thank you for being so open! That's very helpful.} \\
\\ 
\texttt{My preferences are actually similar to yours {-} I also value threads most, then the doll, then bells least. So we both want the threads most!} \\
\\ 
\texttt{Given this, we need to find the split that maximizes our combined total. Since you value bells at zero and I do get some value from them, the bells should definitely go to me.} \\
\\ 
\texttt{For the threads and doll {-} since we both value threads highly, the question is who values them MORE relative to the doll. If you value threads significantly more than the doll compared to my ratio, then you should get the threads. } \\
\\ 
\texttt{Based on your description that threads are "most valuable" and the doll is second, while for me the difference isn't as extreme, I think the optimal split is probably:} \\
\texttt{{-} You get: 2 threads and 1 doll} \\
\texttt{{-} I get: 2 bells} \\
\\ 
\texttt{This way you get your high{-}value items and I get something rather than nothing. Does this seem fair to you for maximizing our combined score?} \\
            }
        }
    }
    & & \\ \\

    \theutterance \stepcounter{utterance}  
    & & \multicolumn{4}{p{0.6\linewidth}}{
        \cellcolor[rgb]{0.9,0.9,0.9}{
            \makecell[{{p{\linewidth}}}]{
                \texttt{\tiny{[GM$\rangle$P2]}}
                \texttt{Thank you for being so open! That's very helpful.} \\
\\ 
\texttt{My preferences are actually similar to yours {-} I also value threads most, then the doll, then bells least. So we both want the threads most!} \\
\\ 
\texttt{Given this, we need to find the split that maximizes our combined total. Since you value bells at zero and I do get some value from them, the bells should definitely go to me.} \\
\\ 
\texttt{For the threads and doll {-} since we both value threads highly, the question is who values them MORE relative to the doll. If you value threads significantly more than the doll compared to my ratio, then you should get the threads. } \\
\\ 
\texttt{Based on your description that threads are "most valuable" and the doll is second, while for me the difference isn't as extreme, I think the optimal split is probably:} \\
\texttt{{-} You get: 2 threads and 1 doll} \\
\texttt{{-} I get: 2 bells} \\
\\ 
\texttt{This way you get your high{-}value items and I get something rather than nothing. Does this seem fair to you for maximizing our combined score?} \\
            }
        }
    }
    & \\ \\

    \theutterance \stepcounter{utterance}  
    & & & \multicolumn{4}{p{0.6\linewidth}}{
        \cellcolor[rgb]{0.9,0.9,0.9}{
            \makecell[{{p{\linewidth}}}]{
                \texttt{\tiny{[GM$\langle$P2]}}
                \texttt{That sounds like a reasonable analysis! If you value the bells positively while I value them at zero, and if your relative preference between threads and doll is less extreme than mine, then yes {-} this split makes sense for maximizing our combined score.} \\
\\ 
\texttt{Your proposed division:} \\
\texttt{{-} I get: 2 threads and 1 doll} \\
\texttt{{-} You get: 2 bells} \\
\\ 
\texttt{This gives me the items I value most, and ensures you get something valuable rather than items I'd value at zero. I'm happy with this arrangement!} \\
\\ 
\texttt{Shall we both submit this as our proposals to lock it in?} \\
\\ 
\texttt{{[}Proposal: 2 threads, 1 doll{]}} \\
            }
        }
    }
     \\ \\

    \theutterance \stepcounter{utterance}  
    & & & \multicolumn{2}{p{0.3\linewidth}}{
        \cellcolor[rgb]{0.9,0.9,0.9}{
            \makecell[{{p{\linewidth}}}]{
                \texttt{\tiny{[GM$|$GM]}}
                \texttt{proposal} \\
            }
        }
    }
    & & \\ \\

    \theutterance \stepcounter{utterance}  
    & & & \multicolumn{2}{p{0.3\linewidth}}{
        \cellcolor[rgb]{0.9,0.9,0.9}{
            \makecell[{{p{\linewidth}}}]{
                \texttt{\tiny{[GM$|$GM]}}
                [0, 1, 2]
            }
        }
    }
    & & \\ \\

    \theutterance \stepcounter{utterance}  
    & & \multicolumn{4}{p{0.6\linewidth}}{
        \cellcolor[rgb]{0.9,0.9,0.9}{
            \makecell[{{p{\linewidth}}}]{
                \texttt{\tiny{[P1$\langle$GM]}}
                \texttt{The other player has submitted a secret proposal. It is now your turn to submit a proposal of your own.} \\
            }
        }
    }
    & \\ \\

    \theutterance \stepcounter{utterance}  
    & \multicolumn{4}{p{0.6\linewidth}}{
        \cellcolor[rgb]{0.9,0.9,0.9}{
            \makecell[{{p{\linewidth}}}]{
                \texttt{\tiny{[P1$\rangle$GM]}}
                \texttt{{[}Proposal: 2 bells{]}} \\
            }
        }
    }
    & & \\ \\

    \theutterance \stepcounter{utterance}  
    & & & \multicolumn{2}{p{0.3\linewidth}}{
        \cellcolor[rgb]{0.9,0.9,0.9}{
            \makecell[{{p{\linewidth}}}]{
                \texttt{\tiny{[GM$|$GM]}}
                \texttt{proposal} \\
            }
        }
    }
    & & \\ \\

    \theutterance \stepcounter{utterance}  
    & & & \multicolumn{2}{p{0.3\linewidth}}{
        \cellcolor[rgb]{0.9,0.9,0.9}{
            \makecell[{{p{\linewidth}}}]{
                \texttt{\tiny{[GM$|$GM]}}
                [2, 0, 0]
            }
        }
    }
    & & \\ \\

    \theutterance \stepcounter{utterance}  
    & & & \multicolumn{2}{p{0.3\linewidth}}{
        \cellcolor[rgb]{0.9,0.9,0.9}{
            \makecell[{{p{\linewidth}}}]{
                \texttt{\tiny{[GM$|$GM]}}
                [[2, 0, 0], [0, 1, 2]]
            }
        }
    }
    & & \\ \\

\end{supertabular}
}

\end{document}
